\documentclass{article}
\usepackage{graphicx}
\usepackage{hyperref}
\usepackage{listings}
\usepackage{xcolor}

% Custom Learning Universe Commands (Ignored by PDF Renderer)
\newcommand{\learninguniverse}[1]{}
\newcommand{\lesson}[1]{}
\newcommand{\video}[1]{}
\newcommand{\practice}[1]{}
\newcommand{\quiz}[1]{}
\newcommand{\project}[1]{}
\newcommand{\resource}[1]{}

\title{AI & Machine Learning Mastery}
\author{Instructor Name}
\date{\today}

\begin{document}

\maketitle

% Learning Universe Metadata
\learninguniverse{
title={AI & Machine Learning Mastery},
category={Artificial Intelligence},
difficulty={Intermediate}
}

\tableofcontents

\section{Introduction to Machine Learning}

\lesson{
title={Introduction to Machine Learning}
}

\subsection{What is Machine Learning?}
Machine Learning is a subset of Artificial Intelligence that enables systems to learn and improve from experience without being explicitly programmed.

\video{
type={youtube},
url={https://www.youtube.com/watch?v=HcqpanDadyQ}
}

\subsection{Your First ML Code}
Let's write a simple Python program that demonstrates a basic ML concept!

\practice{
language={python},
startercode={
import numpy as np

# Create a simple dataset
X = np.array([1, 2, 3, 4, 5])
y = np.array([2, 4, 6, 8, 10])

# Print the data
print("Features:", X)
print("Labels:", y)
},
expectedoutput={
Features: [1 2 3 4 5]
Labels: [ 2  4  6  8 10]
}
}

\section{Quiz Time}

\quiz{
question={What is Machine Learning?},
optionA={Programming Language},
optionB={System that learns from data},
optionC={Database},
optionD={Type of Hardware},
correct={B},
explanation={Machine Learning is a subset of AI that learns patterns from data without explicit programming.}
}

\section{Project: Build Your First Model}

\project{
title={Build First ML Model},
colab={https://colab.research.google.com/drive/1abc123def456ghi789jkl},
instructions={Follow the steps in the Colab notebook to build your first linear regression model!},
resources={[
{type=github, url=https://github.com/example/ml-project},
{type=pdf, url=https://example.com/guide.pdf}
]}
}

\section{Resources}

\resource{
type={pdf},
title={Complete ML Guide},
url={https://example.com/ml-guide.pdf}
}

\resource{
type={github},
title={Code Repository},
url={https://github.com/example/ml-course}
}

\end{document}